\iftestbuild % TEST-ONLY-BEGIN
\setcounter{section}{3}\setcounter{theorem}{58}
\fi % TEST-ONLY-END
% Source: book pp. 82--95 (PDF 96--109); solutions pp. 93--103, 170--174.
\section{Invertibility and Isomorphisms}

\subsection{Invertible Linear Maps}

We begin this section by defining the notions of invertible and inverse in the
context of linear maps.

\begin{definition}[Invertible, inverse]\label{ax:3.59}\leavevmode
\begin{itemize}
\item A linear map $T\in\Lin(V,W)$ is called \emph{invertible} if there exists
  a linear map $S\in\Lin(W,V)$ such that $ST$ equals the identity operator on
  $V$ and $TS$ equals the identity operator on $W$.
\item A linear map $S\in\Lin(W,V)$ satisfying $ST=I$ and $TS=I$ is called an
  \emph{inverse} of $T$ (note that the first $I$ is the identity operator on $V$
  and the second $I$ is the identity operator on $W$).
\end{itemize}
\end{definition}

The definition above mentions ``an inverse''. However, the next result shows
that we can change this terminology to ``the inverse''.

\begin{result}[Inverse is unique]\label{ax:3.60}
An invertible linear map has a unique inverse.
\end{result}

\begin{proof}
Suppose $T\in\Lin(V,W)$ is invertible and $S_1$ and $S_2$ are inverses of $T$.
Then
\[ S_1 = S_1I = S_1(TS_2) = (S_1T)S_2 = IS_2 = S_2. \]
Thus $S_1=S_2$.
\end{proof}

Now that we know that the inverse is unique, we can give it a notation.

\begin{notation}[$T^{-1}$]\label{ax:3.61}
If $T$ is invertible, then its inverse is denoted by $T^{-1}$. In other words,
if $T\in\Lin(V,W)$ is invertible, then $T^{-1}$ is the unique element of
$\Lin(W,V)$ such that $T^{-1}T=I$ and $TT^{-1}=I$.
\end{notation}

\begin{example}[Inverse of a linear map from $\R^3$ to $\R^3$]\label{ax:3.62}
Suppose $T\in\Lin(\R^3)$ is defined by $T(x,y,z)=(-y,x,4z)$. Thus $T$ is a
counterclockwise rotation by $90^\circ$ in the $xy$-plane and a stretch by a
factor of $4$ in the direction of the $z$-axis.

Hence the inverse map $T^{-1}\in\Lin(\R^3)$ is the clockwise rotation by
$90^\circ$ in the $xy$-plane and a stretch by a factor of $1/4$ in the
direction of the $z$-axis:
\[ T^{-1}(x,y,z) = (y,-x,1/4z). \]
\end{example}

The next result shows that a linear map is invertible if and only if it is
one-to-one and onto.

\begin{result}[Invertibility $\iff$ injectivity and surjectivity]\label{ax:3.63}
A linear map is invertible if and only if it is injective and surjective.
\end{result}

\begin{proof}
Suppose $T\in\Lin(V,W)$. We need to show that $T$ is invertible if and only if
it is injective and surjective.

First suppose $T$ is invertible. To show that $T$ is injective, suppose
$u,v\in V$ and $Tu=Tv$. Then
\[ u = T^{-1}(Tu) = T^{-1}(Tv) = v, \]
so $u=v$. Hence $T$ is injective.

We are still assuming that $T$ is invertible. Now we want to prove that $T$ is
surjective. To do this, let $w\in W$. Then $w=T(T^{-1}w)$, which shows that $w$
is in the range of $T$. Thus $\range T=W$. Hence $T$ is surjective, completing
this direction of the proof.

Now suppose $T$ is injective and surjective. We want to prove that $T$ is
invertible. For each $w\in W$, define $S(w)$ to be the unique element of $V$
such that $T\bigl(S(w)\bigr)=w$ (the existence and uniqueness of such an
element follow from the surjectivity and injectivity of $T$). The definition of
$S$ implies that $T\circ S$ equals the identity operator on $W$.

To prove that $S\circ T$ equals the identity operator on $V$, let $v\in V$.
Then
\[ T\bigl((S\circ T)v\bigr) = (T\circ S)(Tv) = I(Tv) = Tv. \]
This equation implies that $(S\circ T)v=v$ (because $T$ is injective). Thus
$S\circ T$ equals the identity operator on $V$.

To complete the proof, we need to show that $S$ is linear. To do this, suppose
$w_1,w_2\in W$. Then
\[ T\bigl(S(w_1)+S(w_2)\bigr) = T\bigl(S(w_1)\bigr)+T\bigl(S(w_2)\bigr)
   = w_1+w_2. \]
Thus $S(w_1)+S(w_2)$ is the unique element of $V$ that $T$ maps to $w_1+w_2$.
By the definition of $S$, this implies that $S(w_1+w_2)=S(w_1)+S(w_2)$. Hence
$S$ satisfies the additive property required for linearity.

The proof of homogeneity is similar. Specifically, if $w\in W$ and
$\lambda\in\F$, then
\[ T\bigl(\lambda S(w)\bigr) = \lambda T\bigl(S(w)\bigr) = \lambda w. \]
Thus $\lambda S(w)$ is the unique element of $V$ that $T$ maps to $\lambda w$.
By the definition of $S$, this implies that $S(\lambda w)=\lambda S(w)$. Hence
$S$ is linear, as desired.
\end{proof}

For a linear map from a vector space to itself, you might wonder whether
injectivity alone, or surjectivity alone, is enough to imply invertibility. On
infinite-dimensional vector spaces, neither condition alone implies
invertibility, as illustrated by the next example, which uses two familiar
linear maps from Example~\ref{ax:3.3}.

\begin{example}[Neither injectivity nor surjectivity implies
invertibility]\label{ax:3.64}\leavevmode
\begin{itemize}
\item The multiplication by $x^2$ linear map from $\Poly(\R)$ to $\Poly(\R)$
  (see \ref{ax:3.3}) is injective but it is not invertible because it is not
  surjective (the polynomial $1$ is not in the range).
\item The backward shift linear map from $\F^\infty$ to $\F^\infty$ (see
  \ref{ax:3.3}) is surjective but it is not invertible because it is not
  injective [the vector $(1,0,0,0,\dots)$ is in the null space].
\end{itemize}
\end{example}

In view of the example above, the next result is remarkable---it states that
for a linear map from a finite-dimensional vector space to a vector space of the
same dimension, either injectivity or surjectivity alone implies the other
condition. Note that the hypothesis below that $\dim V=\dim W$ is automatically
satisfied in the important special case where $V$ is finite-dimensional and
$W=V$.

\begin{result}[Injectivity is equivalent to surjectivity (if
$\dim V=\dim W<\infty$)]\label{ax:3.65}
Suppose that $V$ and $W$ are finite-dimensional vector spaces, $\dim V=\dim W$,
and $T\in\Lin(V,W)$. Then
\[ T \text{ is invertible} \iff T \text{ is injective} \iff
   T \text{ is surjective}. \]
\end{result}

\begin{proof}
The fundamental theorem of linear maps (\ref{ax:3.21}) states that
\[ \refstepcounter{theorem}\tag*{\thetheorem}\label{ax:3.66}
   \dim V = \dim\nul T+\dim\range T. \]
If $T$ is injective (which by \ref{ax:3.15} is equivalent to the condition
$\dim\nul T=0$), then the equation above implies that
\[ \dim\range T = \dim V-\dim\nul T = \dim V = \dim W, \]
which implies that $T$ is surjective (by \ref{ax:2.39}).

Conversely, if $T$ is surjective, then \ref{ax:3.66} implies that
\[ \dim\nul T = \dim V-\dim\range T = \dim V-\dim W = 0, \]
which implies that $T$ is injective.

Thus we have shown that $T$ is injective if and only if $T$ is surjective. Thus
if $T$ is either injective or surjective, then $T$ is both injective and
surjective, which implies that $T$ is invertible. Hence $T$ is invertible if and
only if $T$ is injective if and only if $T$ is surjective.
\end{proof}

The next example illustrates the power of the previous result. Although it is
possible to prove the result in the example below without using linear algebra,
the proof using linear algebra is cleaner and easier.

\begin{example}[There exists a polynomial $p$ such that
$\bigl((x^2+5x+7)p\bigr)''=q$]\label{ax:3.67}
The linear map
\[ p \mapsto \bigl((x^2+5x+7)p\bigr)'' \]
from $\Poly(\R)$ to itself is injective, as you can show. Thus we are tempted
to use \ref{ax:3.65} to show that this map is surjective. However,
Example~\ref{ax:3.64} shows that the magic of \ref{ax:3.65} does not apply to
the infinite-dimensional vector space $\Poly(\R)$. We will get around this
problem by restricting attention to the finite-dimensional vector space
$\Poly_m(\R)$.

Suppose $q\in\Poly(\R)$. There exists a nonnegative integer $m$ such that
$q\in\Poly_m(\R)$. Define $T\colon\Poly_m(\R)\to\Poly_m(\R)$ by
\[ Tp = \bigl((x^2+5x+7)p\bigr)''. \]
Multiplying a nonzero polynomial by $(x^2+5x+7)$ increases the degree by $2$,
and then differentiating twice reduces the degree by $2$. Thus $T$ is indeed a
linear map from $\Poly_m(\R)$ to itself.

Every polynomial whose second derivative equals $0$ is of the form $ax+b$,
where $a,b\in\R$. Thus $\nul T=\{0\}$. Hence $T$ is injective.

Thus $T$ is surjective (by \ref{ax:3.65}), which means that there exists a
polynomial $p\in\Poly_m(\R)$ such that $\bigl((x^2+5x+7)p\bigr)''=q$, as
claimed in the title of this example.
\end{example}

Exercise~35 in Section~6A gives a similar but more spectacular example of using
\ref{ax:3.65}.

The hypothesis in the result below that $\dim V=\dim W$ holds in the important
special case in which $V$ is finite-dimensional and $W=V$. Thus in that case,
the equation $ST=I$ implies that $ST=TS$, even though we do not have
multiplicative commutativity of arbitrary linear maps from $V$ to $V$.

\begin{result}[$ST=I \iff TS=I$ (on vector spaces of the same
dimension)]\label{ax:3.68}
Suppose $V$ and $W$ are finite-dimensional vector spaces of the same dimension,
$S\in\Lin(W,V)$, and $T\in\Lin(V,W)$. Then $ST=I$ if and only if $TS=I$.
\end{result}

\begin{proof}
First suppose $ST=I$. If $v\in V$ and $Tv=0$, then
\[ v = Iv = (ST)v = S(Tv) = S(0) = 0. \]
Thus $T$ is injective (by \ref{ax:3.15}). Because $V$ and $W$ have the same
dimension, this implies that $T$ is invertible (by \ref{ax:3.65}).

Now multiply both sides of the equation $ST=I$ by $T^{-1}$ on the right,
getting
\[ S = T^{-1}. \]
Thus $TS=TT^{-1}=I$, as desired.

To prove the implication in the other direction, simply reverse the roles of
$S$ and $T$ (and $V$ and $W$) in the direction we have already proved, showing
that if $TS=I$, then $ST=I$.
\end{proof}

\subsection{Isomorphic Vector Spaces}

The next definition captures the idea of two vector spaces that are essentially
the same, except for the names of their elements.

\begin{definition}[Isomorphism, isomorphic]\label{ax:3.69}\leavevmode
\begin{itemize}
\item An \emph{isomorphism} is an invertible linear map.
\item Two vector spaces are called \emph{isomorphic} if there is an isomorphism
  from one vector space onto the other one.
\end{itemize}
\end{definition}

Think of an isomorphism $T\colon V\to W$ as relabeling $v\in V$ as $Tv\in W$.
This viewpoint explains why two isomorphic vector spaces have the same vector
space properties. The terms ``isomorphism'' and ``invertible linear map'' mean
the same thing. Use ``isomorphism'' when you want to emphasize that the two
spaces are essentially the same.

It can be difficult to determine whether two mathematical structures (such as
groups or topological spaces) are essentially the same, differing only in the
names of the elements of underlying sets. However, the next result shows that
we need to look at only a single number (the dimension) to determine whether two
vector spaces are isomorphic.

\begin{result}[Dimension shows whether vector spaces are
isomorphic]\label{ax:3.70}
Two finite-dimensional vector spaces over $\F$ are isomorphic if and only if
they have the same dimension.
\end{result}

\begin{proof}
First suppose $V$ and $W$ are isomorphic finite-dimensional vector spaces. Thus
there exists an isomorphism $T$ from $V$ onto $W$. Because $T$ is invertible,
we have $\nul T=\{0\}$ and $\range T=W$. Thus
\[ \dim\nul T = 0 \quad\text{and}\quad \dim\range T = \dim W. \]
The formula
\[ \dim V = \dim\nul T+\dim\range T \]
(the fundamental theorem of linear maps, which is \ref{ax:3.21}) thus becomes
the equation $\dim V=\dim W$, completing the proof in one direction.

To prove the other direction, suppose $V$ and $W$ are finite-dimensional vector
spaces of the same dimension. Let $v_1,\dots,v_n$ be a basis of $V$ and
$w_1,\dots,w_n$ be a basis of $W$. Let $T\in\Lin(V,W)$ be defined by
\[ T(c_1v_1+\cdots+c_nv_n) = c_1w_1+\cdots+c_nw_n. \]
Then $T$ is a well-defined linear map because $v_1,\dots,v_n$ is a basis of
$V$. Also, $T$ is surjective because $w_1,\dots,w_n$ spans $W$. Furthermore,
$\nul T=\{0\}$ because $w_1,\dots,w_n$ is linearly independent. Thus $T$ is
injective. Because $T$ is injective and surjective, it is an isomorphism (see
\ref{ax:3.63}). Hence $V$ and $W$ are isomorphic.
\end{proof}

\marginnote[-15mm]{Every finite-dimensional vector space is isomorphic to some $\F^n$.
Thus why not just study $\F^n$ instead of more general vector spaces? To answer
this question, note that an investigation of $\F^n$ would soon lead to other
vector spaces. For example, we would encounter the null space and range of
linear maps. Although each of these vector spaces is isomorphic to some $\F^m$,
thinking of them that way often adds complexity but no new insight.}%
The previous result implies that each finite-dimensional vector space $V$ is
isomorphic to $\F^n$, where $n=\dim V$. For example, if $m$ is a nonnegative
integer, then $\Poly_m(\F)$ is isomorphic to $\F^{m+1}$.

Recall that the notation $\F^{m,n}$ denotes the vector space of $m$-by-$n$
matrices with entries in $\F$. If $v_1,\dots,v_n$ is a basis of $V$ and
$w_1,\dots,w_m$ is a basis of $W$, then for each $T\in\Lin(V,W)$, we have a
matrix $\Mat(T)\in\F^{m,n}$. Thus once bases have been fixed for $V$ and $W$,
$\Mat$ becomes a function from $\Lin(V,W)$ to $\F^{m,n}$. Notice that
\ref{ax:3.35} and \ref{ax:3.38} show that $\Mat$ is a linear map. This linear
map is actually an isomorphism, as we now show.

\begin{result}[$\Lin(V,W)$ and $\F^{m,n}$ are isomorphic]\label{ax:3.71}
Suppose $v_1,\dots,v_n$ is a basis of $V$ and $w_1,\dots,w_m$ is a basis of
$W$. Then $\Mat$ is an isomorphism between $\Lin(V,W)$ and $\F^{m,n}$.
\end{result}

\begin{proof}
We already noted that $\Mat$ is linear. We need to prove that $\Mat$ is
injective and surjective.

We begin with injectivity. If $T\in\Lin(V,W)$ and $\Mat(T)=0$, then $Tv_k=0$
for each $k=1,\dots,n$. Because $v_1,\dots,v_n$ is a basis of $V$, this implies
$T=0$. Thus $\Mat$ is injective (by \ref{ax:3.15}).

To prove that $\Mat$ is surjective, suppose $A\in\F^{m,n}$. By the linear map
lemma (\ref{ax:3.4}), there exists $T\in\Lin(V,W)$ such that
\[ Tv_k = \sum_{j=1}^m A_{j,k}w_j \]
for each $k=1,\dots,n$. Because $\Mat(T)$ equals $A$, the range of $\Mat$
equals $\F^{m,n}$, as desired.
\end{proof}

Now we can determine the dimension of the vector space of linear maps from one
finite-dimensional vector space to another.

\begin{result}[$\dim\Lin(V,W)=(\dim V)(\dim W)$]\label{ax:3.72}
Suppose $V$ and $W$ are finite-dimensional. Then $\Lin(V,W)$ is
finite-dimensional and
\[ \dim\Lin(V,W) = (\dim V)(\dim W). \]
\end{result}

\begin{proof}
The desired result follows from \ref{ax:3.71}, \ref{ax:3.70}, and
\ref{ax:3.40}.
\end{proof}

\subsection{Linear Maps Thought of as Matrix Multiplication}

Previously we defined the matrix of a linear map. Now we define the matrix of a
vector.

\begin{definition}[Matrix of a vector, $\Mat(v)$]\label{ax:3.73}
Suppose $v\in V$ and $v_1,\dots,v_n$ is a basis of $V$. The \emph{matrix of $v$}
with respect to this basis is the $n$-by-$1$ matrix
\[ \Mat(v) = \begin{pmatrix} b_1\\ \vdots\\ b_n \end{pmatrix}, \]
where $b_1,\dots,b_n$ are the scalars such that
\[ v = b_1v_1+\cdots+b_nv_n. \]
\end{definition}

The matrix $\Mat(v)$ of a vector $v\in V$ depends on the basis $v_1,\dots,v_n$
of $V$, as well as on $v$. However, the basis should be clear from the context
and thus it is not included in the notation.

\begin{example}[Matrix of a vector]\label{ax:3.74}\leavevmode
\begin{itemize}
\item The matrix of the polynomial $2-7x+5x^3+x^4$ with respect to the standard
  basis of $\Poly_4(\R)$ is
  \[ \begin{pmatrix} 2\\ -7\\ 0\\ 5\\ 1 \end{pmatrix}. \]
\item The matrix of a vector $x\in\F^n$ with respect to the standard basis is
  obtained by writing the coordinates of $x$ as the entries in an $n$-by-$1$
  matrix. In other words, if $x=(x_1,\dots,x_n)\in\F^n$, then
  \[ \Mat(x) = \begin{pmatrix} x_1\\ \vdots\\ x_n \end{pmatrix}. \]
\end{itemize}
\end{example}

Occasionally we want to think of elements of $V$ as relabeled to be $n$-by-$1$
matrices. Once a basis $v_1,\dots,v_n$ is chosen, the function $\Mat$ that
takes $v\in V$ to $\Mat(v)$ is an isomorphism of $V$ onto $\F^{n,1}$ that
implements this relabeling.

Recall that if $A$ is an $m$-by-$n$ matrix, then $A_{\cdot,k}$ denotes the
$k^{\text{th}}$ column of $A$, thought of as an $m$-by-$1$ matrix. In the next
result, $\Mat(Tv_k)$ is computed with respect to the basis $w_1,\dots,w_m$ of
$W$.

\begin{result}[$\Mat(T)_{\cdot,k}=\Mat(Tv_k)$]\label{ax:3.75}
Suppose $T\in\Lin(V,W)$ and $v_1,\dots,v_n$ is a basis of $V$ and
$w_1,\dots,w_m$ is a basis of $W$. Let $1\le k\le n$. Then the
$k^{\text{th}}$ column of $\Mat(T)$, which is denoted by $\Mat(T)_{\cdot,k}$,
equals $\Mat(Tv_k)$.
\end{result}

\begin{proof}
The desired result follows immediately from the definitions of $\Mat(T)$ and
$\Mat(Tv_k)$.
\end{proof}

The next result shows how the notions of the matrix of a linear map, the matrix
of a vector, and matrix multiplication fit together.

\begin{result}[Linear maps act like matrix multiplication]\label{ax:3.76}
Suppose $T\in\Lin(V,W)$ and $v\in V$. Suppose $v_1,\dots,v_n$ is a basis of $V$
and $w_1,\dots,w_m$ is a basis of $W$. Then
\[ \Mat(Tv) = \Mat(T)\Mat(v). \]
\end{result}

\begin{proof}
Suppose $v=b_1v_1+\cdots+b_nv_n$, where $b_1,\dots,b_n\in\F$. Thus
\[ \refstepcounter{theorem}\tag*{\thetheorem}\label{ax:3.77}
   Tv = b_1Tv_1+\cdots+b_nTv_n. \]
Hence
\begin{align*}
\Mat(Tv) &= b_1\Mat(Tv_1)+\cdots+b_n\Mat(Tv_n)\\
         &= b_1\Mat(T)_{\cdot,1}+\cdots+b_n\Mat(T)_{\cdot,n}\\
         &= \Mat(T)\Mat(v),
\end{align*}
where the first equality follows from \ref{ax:3.77} and the linearity of
$\Mat$, the second equality comes from \ref{ax:3.75}, and the last equality
comes from \ref{ax:3.50}.
\end{proof}

Each $m$-by-$n$ matrix $A$ induces a linear map from $\F^{n,1}$ to $\F^{m,1}$,
namely the matrix multiplication function that takes $x\in\F^{n,1}$ to
$Ax\in\F^{m,1}$. The result above can be used to think of every linear map
(from a finite-dimensional vector space to another finite-dimensional vector
space) as a matrix multiplication map after suitable relabeling via the
isomorphisms given by $\Mat$. Specifically, if $T\in\Lin(V,W)$ and we identify
$v\in V$ with $\Mat(v)\in\F^{n,1}$, then the result above says that we can
identify $Tv$ with $\Mat(T)\Mat(v)$.

Because the result above allows us to think (via isomorphisms) of each linear
map as multiplication on $\F^{n,1}$ by some matrix $A$, keep in mind that the
specific matrix $A$ depends not only on the linear map but also on the choice
of bases. One of the themes of many of the most important results in later
chapters will be the choice of a basis that makes the matrix $A$ as simple as
possible.

In this book, we concentrate on linear maps rather than on matrices. However,
sometimes thinking of linear maps as matrices (or thinking of matrices as
linear maps) gives important insights that we will find useful.

Notice that no bases are in sight in the statement of the next result. Although
$\Mat(T)$ in the next result depends on a choice of bases of $V$ and $W$, the
next result shows that the column rank of $\Mat(T)$ is the same for all such
choices (because $\range T$ does not depend on a choice of basis).

\begin{result}[Dimension of $\range T$ equals column rank of
$\Mat(T)$]\label{ax:3.78}
Suppose $V$ and $W$ are finite-dimensional and $T\in\Lin(V,W)$. Then
$\dim\range T$ equals the column rank of $\Mat(T)$.
\end{result}

\begin{proof}\emergencystretch=1.5em
Suppose $v_1,\dots,v_n$ is a basis of $V$ and $w_1,\dots,w_m$ is a basis of
$W$. The linear map that takes $w\in W$ to $\Mat(w)$ is an isomorphism from $W$
onto the space $\F^{m,1}$ of $m$-by-$1$ column vectors. The restriction of this
isomorphism to $\range T$ [which equals $\spn(Tv_1,\dots,Tv_n)$ by Exercise~10
in Section~3B] is an isomorphism from $\range T$ onto
$\spn\bigl(\Mat(Tv_1),\dots,\Mat(Tv_n)\bigr)$. For each $k\in\{1,\dots,n\}$,
the $m$-by-$1$ matrix $\Mat(Tv_k)$ equals column $k$ of $\Mat(T)$. Thus
\[ \dim\range T = \text{the column rank of } \Mat(T), \]
as desired.
\end{proof}

\subsection{Change of Basis}

In Section~3C we defined the matrix
\[ \Mat\bigl(T,(v_1,\dots,v_n),(w_1,\dots,w_m)\bigr) \]
of a linear map $T$ from $V$ to a possibly different vector space $W$, where
$v_1,\dots,v_n$ is a basis of $V$ and $w_1,\dots,w_m$ is a basis of $W$. For
linear maps from a vector space to itself, we usually use the same basis for
both the domain vector space and the target vector space. When using a single
basis in both capacities, we often write the basis only once. In other words,
if $T\in\Lin(V)$ and $v_1,\dots,v_n$ is a basis of $V$, then the notation
$\Mat\bigl(T,(v_1,\dots,v_n)\bigr)$ is defined by the equation
\[ \Mat\bigl(T,(v_1,\dots,v_n)\bigr)
   = \Mat\bigl(T,(v_1,\dots,v_n),(v_1,\dots,v_n)\bigr). \]
If the basis $v_1,\dots,v_n$ is clear from the context, then we can write just
$\Mat(T)$.

\begin{definition}[Identity matrix, $I$]\label{ax:3.79}
Suppose $n$ is a positive integer. The $n$-by-$n$ matrix
\[ \begin{pmatrix} 1 & & 0\\ & \ddots & \\ 0 & & 1 \end{pmatrix} \]
with $1$'s on the diagonal (the entries where the row number equals the column
number) and $0$'s elsewhere is called the \emph{identity matrix} and is denoted
by $I$.
\end{definition}

In the definition above, the $0$ in the lower left corner of the matrix
indicates that all entries below the diagonal are $0$, and the $0$ in the upper
right corner indicates that all entries above the diagonal are $0$.

With respect to each basis of $V$, the matrix of the identity operator
$I\in\Lin(V)$ is the identity matrix $I$. Note that the symbol $I$ is used to
denote both the identity operator and the identity matrix. The context indicates
which meaning of $I$ is intended. For example, consider the equation
$\Mat(I)=I$; on the left side $I$ denotes the identity operator, and on the
right side $I$ denotes the identity matrix.

If $A$ is a square matrix (meaning it has the same number of rows as columns)
with the same size as $I$, then $AI=IA=A$, as you should verify.

\begin{definition}[Invertible, inverse, $A^{-1}$]\label{ax:3.80}
A square matrix $A$ is called \emph{invertible} if there is a square matrix $B$
of the same size such that $AB=BA=I$; we call $B$ the \emph{inverse} of $A$ and
denote it by $A^{-1}$.
\end{definition}

\marginnote{Some mathematicians use the terms \textbf{nonsingular} and
\textbf{singular}, which mean the same as invertible and noninvertible.}%
The same proof as used in \ref{ax:3.60} shows that if $A$ is an invertible
square matrix, then there is a unique matrix $B$ such that $AB=BA=I$ (and thus
the notation $B=A^{-1}$ is justified).

If $A$ is an invertible matrix, then $(A^{-1})^{-1}=A$ because
\[ A^{-1}A = AA^{-1} = I. \]
Also, if $A$ and $C$ are invertible square matrices of the same size, then $AC$
is invertible and $(AC)^{-1}=C^{-1}A^{-1}$ because
\begin{align*}
(AC)(C^{-1}A^{-1}) &= A(CC^{-1})A^{-1}\\
                   &= AIA^{-1}\\
                   &= AA^{-1}\\
                   &= I,
\end{align*}
and similarly $(C^{-1}A^{-1})(AC)=I$.

The next result holds because we defined matrix multiplication to make it
true---see \ref{ax:3.43} and the material preceding it. Now we are just being
more explicit about the bases involved.

\begin{result}[Matrix of product of linear maps]\label{ax:3.81}
Suppose $T\in\Lin(U,V)$ and $S\in\Lin(V,W)$. If $u_1,\dots,u_m$ is a basis of
$U$, $v_1,\dots,v_n$ is a basis of $V$, and $w_1,\dots,w_p$ is a basis of $W$,
then
\begin{multline*}
\Mat\bigl(ST,(u_1,\dots,u_m),(w_1,\dots,w_p)\bigr) =\\
\Mat\bigl(S,(v_1,\dots,v_n),(w_1,\dots,w_p)\bigr)
\Mat\bigl(T,(u_1,\dots,u_m),(v_1,\dots,v_n)\bigr).
\end{multline*}
\end{result}

The next result deals with the matrix of the identity operator $I$ with respect
to two different bases. Note that the $k^{\text{th}}$ column of
$\Mat\bigl(I,(u_1,\dots,u_n),(v_1,\dots,v_n)\bigr)$ consists of the scalars
needed to write $u_k$ as a linear combination of the basis $v_1,\dots,v_n$.

In the statement of the next result, $I$ denotes the identity operator from $V$
to $V$. In the proof, $I$ also denotes the $n$-by-$n$ identity matrix.

\begin{result}[Matrix of identity operator with respect to two
bases]\label{ax:3.82}
Suppose that $u_1,\dots,u_n$ and $v_1,\dots,v_n$ are bases of $V$. Then the
matrices
\[ \Mat\bigl(I,(u_1,\dots,u_n),(v_1,\dots,v_n)\bigr) \quad\text{and}\quad
   \Mat\bigl(I,(v_1,\dots,v_n),(u_1,\dots,u_n)\bigr) \]
are invertible, and each is the inverse of the other.
\end{result}

\begin{proof}
In \ref{ax:3.81}, replace $w_k$ with $u_k$, and replace $S$ and $T$ with $I$,
getting
\[ I = \Mat\bigl(I,(v_1,\dots,v_n),(u_1,\dots,u_n)\bigr)
       \Mat\bigl(I,(u_1,\dots,u_n),(v_1,\dots,v_n)\bigr). \]
Now interchange the roles of the $u$'s and $v$'s, getting
\[ I = \Mat\bigl(I,(u_1,\dots,u_n),(v_1,\dots,v_n)\bigr)
       \Mat\bigl(I,(v_1,\dots,v_n),(u_1,\dots,u_n)\bigr). \]
These two equations above give the desired result.
\end{proof}

\begin{example}[Matrix of identity operator on $\F^2$ with respect to two
bases]\label{ax:3.83}
Consider the bases $(4,2),(5,3)$ and $(1,0),(0,1)$ of $\F^2$. Because
$I(4,2)=4(1,0)+2(0,1)$ and $I(5,3)=5(1,0)+3(0,1)$, we have
\[ \Mat\Bigl(I,\bigl((4,2),(5,3)\bigr),\bigl((1,0),(0,1)\bigr)\Bigr)
   = \begin{pmatrix} 4 & 5\\ 2 & 3 \end{pmatrix}. \]
The inverse of the matrix above is
\[ \begin{pmatrix} 3/2 & -5/2\\ -1 & 2 \end{pmatrix}, \]
as you should verify. Thus \ref{ax:3.82} implies that
\[ \Mat\Bigl(I,\bigl((1,0),(0,1)\bigr),\bigl((4,2),(5,3)\bigr)\Bigr)
   = \begin{pmatrix} 3/2 & -5/2\\ -1 & 2 \end{pmatrix}. \]
\end{example}

Our next result shows how the matrix of $T$ changes when we change bases. In the
next result, we have two different bases of $V$, each of which is used as a
basis for the domain space and as a basis for the target space. Recall our
shorthand notation that allows us to display a basis only once when it is used
in both capacities:
\[ \Mat\bigl(T,(u_1,\dots,u_n)\bigr)
   = \Mat\bigl(T,(u_1,\dots,u_n),(u_1,\dots,u_n)\bigr). \]

\begin{result}[Change-of-basis formula]\label{ax:3.84}
Suppose $T\in\Lin(V)$. Suppose $u_1,\dots,u_n$ and $v_1,\dots,v_n$ are bases of
$V$. Let
\[ A = \Mat\bigl(T,(u_1,\dots,u_n)\bigr) \quad\text{and}\quad
   B = \Mat\bigl(T,(v_1,\dots,v_n)\bigr) \]
and $C=\Mat\bigl(I,(u_1,\dots,u_n),(v_1,\dots,v_n)\bigr)$. Then
\[ A = C^{-1}BC. \]
\end{result}

\begin{proof}
In \ref{ax:3.81}, replace $w_k$ with $u_k$ and replace $S$ with $I$, getting
\[ \refstepcounter{theorem}\tag*{\thetheorem}\label{ax:3.85}
   A = C^{-1}\Mat\bigl(T,(u_1,\dots,u_n),(v_1,\dots,v_n)\bigr), \]
where we have used \ref{ax:3.82}.

Again use \ref{ax:3.81}, this time replacing $w_k$ with $v_k$. Also replace $T$
with $I$ and replace $S$ with $T$, getting
\[ \Mat\bigl(T,(u_1,\dots,u_n),(v_1,\dots,v_n)\bigr) = BC. \]
Substituting the equation above into \ref{ax:3.85} gives the equation
$A=C^{-1}BC$.
\end{proof}

The proof of the next result is left as an exercise.

\begin{result}[Matrix of inverse equals inverse of matrix]\label{ax:3.86}
Suppose that $v_1,\dots,v_n$ is a basis of $V$ and $T\in\Lin(V)$ is invertible.
Then $\Mat(T^{-1})=\bigl(\Mat(T)\bigr)^{-1}$, where both matrices are with
respect to the basis $v_1,\dots,v_n$.
\end{result}

% ------------------------------------------------------------------ exercises
\begin{exc}

\begin{exercise}
Suppose $T\in\Lin(V,W)$ is invertible. Show that $T^{-1}$ is invertible and
\[ (T^{-1})^{-1} = T. \]
\end{exercise}
\begin{solution}
Since $T$ is invertible, then $T^{-1}\in\Lin(W,V)$. By $TT^{-1}=I$ and
$T^{-1}T=I$, we have that $T^{-1}$ is invertible and $(T^{-1})^{-1}=T$.
\end{solution}

\begin{exercise}
Suppose $T\in\Lin(U,V)$ and $S\in\Lin(V,W)$ are both invertible linear maps.
Prove that $ST\in\Lin(U,W)$ is invertible and that $(ST)^{-1}=T^{-1}S^{-1}$.
\end{exercise}
\begin{solution}
Since $T$ and $S$ are invertible, then $T^{-1}\in\Lin(V,U)$ and
$S^{-1}\in\Lin(W,V)$. By the definition of the inverse, we have that
\[ (ST)(T^{-1}S^{-1}) = S(TT^{-1})S^{-1} = SIS^{-1} = SS^{-1} = I, \]
and
\[ (T^{-1}S^{-1})(ST) = T^{-1}(S^{-1}S)T = T^{-1}IT = T^{-1}T = I. \]
Thus, $ST$ is invertible and $(ST)^{-1}=T^{-1}S^{-1}$.
\end{solution}

\begin{exercise}
Suppose $V$ is finite-dimensional and $T\in\Lin(V)$. Prove that the following
are equivalent.
\begin{parts}
\item $T$ is invertible.
\item $Tv_1,\dots,Tv_n$ is a basis of $V$ for every basis $v_1,\dots,v_n$ of
  $V$.
\item $Tv_1,\dots,Tv_n$ is a basis of $V$ for some basis $v_1,\dots,v_n$ of $V$.
\end{parts}
\end{exercise}
\begin{solution}\leavevmode
\begin{itemize}
\item (a) $\implies$ (b): Suppose $T$ is invertible. Consider scalars
  $\alpha_1,\dots,\alpha_n$ such that
  \[ \alpha_1Tv_1+\cdots+\alpha_nTv_n = 0. \]
  Applying $T^{-1}$ to both sides, we get
  \[ \alpha_1v_1+\cdots+\alpha_nv_n = 0. \]
  Since $v_1,\dots,v_n$ is a basis of $V$, it follows that
  $\alpha_1=\cdots=\alpha_n=0$. Hence, $Tv_1,\dots,Tv_n$ is linearly
  independent. Since $Tv_1,\dots,Tv_n$ has $n$ elements and $V$ is
  $n$-dimensional, it follows that $Tv_1,\dots,Tv_n$ is a basis of $V$.
\item (b) $\implies$ (c): This is trivial.
\item (c) $\implies$ (a): Suppose $Tv_1,\dots,Tv_n$ is a basis of $V$ for some
  basis $v_1,\dots,v_n$ of $V$. Then $\range T=V$. By the fundamental theorem of
  linear maps, we have that $\dim\nul T=0$. Thus, $T$ is injective. Since $T$ is
  injective and $\range T=V$, then $T$ is invertible.
\end{itemize}
\end{solution}

\begin{exercise}
Suppose $V$ is finite-dimensional and $\dim V>1$. Prove that the set of
noninvertible linear maps from $V$ to itself is not a subspace of $\Lin(V)$.
\end{exercise}
\begin{solution}
Let $v_1,\dots,v_n$ be a basis of $V$. Define $T,S\in\Lin(V)$ such that
\[ Tv_1 = 0, \quad Tv_k = v_k \text{ for } k=2,\dots,n, \]
and
\[ Sv_1 = v_1, \quad Sv_k = 0 \text{ for } k=2,\dots,n. \]
Then $T$ and $S$ are non-invertible linear maps, since $\nul T=\spn v_1$ and
$\nul S=\spn(v_2,\dots,v_n)$. However, $T+S$ is invertible, since
\[ (T+S)v_1 = v_1, \quad (T+S)v_k = v_k \text{ for } k=2,\dots,n. \]
Hence, $T+S=I$ is invertible. Therefore, the set of non-invertible linear maps
from $V$ to itself is not a subspace of $\Lin(V)$.
\end{solution}

\begin{exercise}
Suppose $V$ is finite-dimensional, $U$ is a subspace of $V$, and
$S\in\Lin(U,V)$. Prove that there exists an invertible linear map $T$ from $V$
to itself such that $Tu=Su$ for every $u\in U$ if and only if $S$ is injective.
\end{exercise}
\begin{solution}
Suppose there exists an invertible linear map $T$ from $V$ to itself such that
$Tu=Su$ for every $u\in U$. If $Su=0$, then $Tu=0$, and since $T$ is
invertible, $u=0$. Hence, $S$ is injective.

Conversely, suppose $S$ is injective. Let $u_1,\dots,u_m$ be a basis of $U$.
Since $S$ is injective, then $Su_1,\dots,Su_m$ is linearly independent. Extend
$u_1,\dots,u_m$ to a basis $u_1,\dots,u_n$ of $V$. Extend $Su_1,\dots,Su_m$ to
a basis $v_1,\dots,v_n$ of $V$. Define $T\in\Lin(V)$ such that
\[ Tu_k = Su_k \text{ for } k=1,\dots,m, \quad
   Tu_k = v_k \text{ for } k=m+1,\dots,n. \]
Then $T$ is invertible, since $Tu_1,\dots,Tu_n$ is a basis of $V$. Thus, there
exists an invertible linear map $T$ from $V$ to itself such that $Tu=Su$ for
every $u\in U$.
\end{solution}

\begin{exercise}
Suppose that $W$ is finite-dimensional and $S,T\in\Lin(V,W)$. Prove that
$\nul S=\nul T$ if and only if there exists an invertible $E\in\Lin(W)$ such
that $S=ET$.
\end{exercise}
\begin{solution}
Suppose $\nul S=\nul T$. Define a linear map $E_0$ from $\range T$ to
$\range S$ by setting $E_0(Tv)=Sv$ for all $v\in V$. This is well-defined
because if $Tv=Tv'$ for some $v,v'\in V$, then $T(v-v')=0$, so
$v-v'\in\nul T=\nul S$, and thus $S(v-v')=0$, which implies $Sv=Sv'$. Moreover,
if $E_0(Tv)=0$ for some $v\in V$, then $Sv=0$, which implies
$v\in\nul S=\nul T$, and thus $Tv=0$. Hence, $E_0$ is injective. By
Exercise~3D5, we can extend $E_0$ to an invertible linear map $E$ on all of
$W$, and then we have $S=ET$.

Conversely, suppose there exists an invertible $E\in\Lin(W)$ such that $S=ET$.
If $v\in\nul S$, then $Sv=0$, and since $S=ET$, we have $ETv=0$. Since $E$ is
invertible, $Tv=0$ so that $v\in\nul T$. Hence, $\nul S\subseteq\nul T$.
Similarly, if $v\in\nul T$, then $Tv=0$, and since $S=ET$, we have
$Sv=ETv=E0=0$. Thus, $v\in\nul S$, and $\nul T\subseteq\nul S$. Therefore,
$\nul S=\nul T$.
\end{solution}

\begin{exercise}
Suppose that $V$ is finite-dimensional and $S,T\in\Lin(V,W)$. Prove that
$\range S=\range T$ if and only if there exists an invertible $E\in\Lin(V)$ such
that $S=TE$.
\end{exercise}
\begin{solution}
Suppose $\range S=\range T$. Then $\dim\nul S=\dim\nul T$. Let $u_1,\dots,u_k$
be a basis of $\nul S$ and $v_1,\dots,v_k$ be a basis of $\nul T$. Extend both
bases to a basis of $V$, say $u_1,\dots,u_n$ and $v_1,\dots,v_n$. Observe that
for $i>k$, we have that $Su_i\in\range S=\range T$. Then there must exist
$y_i\in V$ such that $Ty_i=Su_i$. Define a linear map $E\in\Lin(V)$ by setting
$Eu_i=v_i$ for $i\le k$ and $Eu_i=y_i$ for $i>k$. To see that $TE=S$, we have
for $i\le k$ that $TEu_i=Tv_i=0=Su_i$. Moreover, for $i>k$, we have
$TEu_i=Ty_i=Su_i$. To see that $E$ is invertible, suppose we have scalars
$\alpha_1,\dots,\alpha_n$ such that
\[ \alpha_1Eu_1+\alpha_2Eu_2+\cdots+\alpha_nEu_n = 0. \]
Applying $T$, we get
\[ \alpha_{k+1}Su_{k+1}+\cdots+\alpha_nSu_n = 0, \]
where we used the fact that $TE=S$ and $Su_1=\cdots=Su_k=0$ because
$u_1,\dots,u_k\in\nul S$. Since $u_1,\dots,u_n$ is a basis of $V$, the vectors
$u_{k+1},\dots,u_n$ map to a spanning list of $\range S$. Since
$\dim\range S=n-k$, it follows that $Su_{k+1},\dots,Su_n$ are linearly
independent. Then from the previous equation, $\alpha_{k+1}=\cdots=\alpha_n=0$.
Substituting this back into the original equation, we get
$\alpha_1v_1+\cdots+\alpha_kv_k=0$, and since $v_1,\dots,v_k$ are linearly
independent, we must have $\alpha_1=\cdots=\alpha_k=0$. Hence, all the scalars
are zero, proving that $E$ is invertible.

Conversely, suppose there exists an invertible $E\in\Lin(V)$ such that $S=TE$.
If $w\in\range S$, then there exists some $v\in V$ such that $Sv=w$. Since
$S=TE$, we have $TEv=w$, and thus $w\in\range T$. Hence,
$\range S\subseteq\range T$. Similarly, if $w\in\range T$, then there exists
some $v\in V$ such that $Tv=w$. Since $E$ is invertible, there exists some
$u\in V$ such that $Eu=v$. Then we have
\[ Su = TEu = Tv = w, \]
and thus $w\in\range S$. Hence, $\range T\subseteq\range S$. Therefore,
$\range S=\range T$.
\end{solution}

\begin{exercise}
Suppose $V$ and $W$ are finite-dimensional and $S,T\in\Lin(V,W)$. Prove that
there exist invertible $E_1\in\Lin(V)$ and $E_2\in\Lin(W)$ such that
$S=E_2TE_1$ if and only if $\dim\nul S=\dim\nul T$.
\end{exercise}
\begin{solution}
Suppose there exist invertible $E_1\in\Lin(V)$ and invertible $E_2\in\Lin(W)$
such that $S=E_2TE_1$. We claim that $E_1$ is an invertible map from $\nul S$
to $\nul T$. To see this, suppose $v,w\in\nul S$ and $E_1v=E_1w$. Then $v=w$
since $E_1$ is invertible, proving injectivity. To see surjectivity, let
$u\in\nul T$. Since $E_1$ is invertible, there exists some $v\in V$ such that
$E_1v=u$. Then
\[ Sv = E_2TE_1v = E_2Tu = E_20 = 0, \]
and thus $v\in\nul S$. Hence, $E_1$ is an invertible map from $\nul S$ to
$\nul T$, and thus $\dim\nul S=\dim\nul T$.

Conversely, suppose $\dim\nul S=\dim\nul T$. Then we may construct an
invertible linear map $E_1\in\Lin(V)$ that maps $\nul S$ onto $\nul T$. We now
claim that $\nul TE_1=\nul S$. To see this, let $v\in\nul TE_1$. This is only
true when $TE_1v=0$, or $E_1v\in\nul T$. By invertibility of $E_1$, we have
$v\in\nul S$. By Exercise~3D6 with $\nul S=\nul TE_1$, there exists an
invertible $E_2\in\Lin(W)$ such that $S=E_2TE_1$.
\end{solution}

\begin{exercise}
Suppose $V$ is finite-dimensional and $T\colon V\to W$ is a surjective linear
map of $V$ onto $W$. Prove that there is a subspace $U$ of $V$ such that
$T|_U$ is an isomorphism of $U$ onto $W$.
\exnote{Here $T|_U$ means the function $T$ restricted to $U$. Thus $T|_U$ is
the function whose domain is $U$, with $T|_U$ defined by $T|_U(u)=Tu$ for every
$u\in U$.}
\end{exercise}
\begin{solution}
Since $T$ is surjective, then $\range T=W$. By Exercise~3B11, there exists a
subspace $U$ of $V$ such that $V=\nul T\oplus U$. We show that $T|_U$ is an
isomorphism of $U$ onto $W$.

Suppose $u\in U$ such that $T|_Uu=0$. Then $Tu=0$, and thus
$u\in\nul T\cap U$. Since $V=\nul T\oplus U$, we have $\nul T\cap U=\{0\}$, and
thus $u=0$. This shows that $T|_U$ is injective.

To see surjectivity, let $w\in W$. Since $T$ is surjective, there exists
$v\in V$ such that $Tv=w$. Write $v=n+u$ with $n\in\nul T$ and $u\in U$. Then
$Tu=T(v-n)=Tv-Tn=w-0=w$. Hence, $T|_U$ is surjective onto $W$, and therefore an
isomorphism of $U$ onto $W$.
\end{solution}

\begin{exercise}
Suppose $V$ and $W$ are finite-dimensional and $U$ is a subspace of $V$. Let
\[ \mathcal{E} = \{T\in\Lin(V,W) : U\subseteq\nul T\}. \]
\begin{parts}
\item Show that $\mathcal{E}$ is a subspace of $\Lin(V,W)$.
\item Find a formula for $\dim\mathcal{E}$ in terms of $\dim V$, $\dim W$, and
  $\dim U$.
\end{parts}
\exnote{Hint: Define $\Phi\colon\Lin(V,W)\to\Lin(U,W)$ by $\Phi(T)=T|_U$. What
is $\nul\Phi$? What is $\range\Phi$?}
\end{exercise}
\begin{solution}\leavevmode
\begin{parts}
\item Consider the zero map $0\in\Lin(V,W)$. For this map, $U\subseteq\nul 0$
  since $\nul 0=V$. Hence, $0\in\mathcal{E}$, showing that $\mathcal{E}$ is
  non-empty. To see that $\mathcal{E}$ is closed under addition and scalar
  multiplication, let $S,T\in\mathcal{E}$ and $\alpha\in\F$. For any $u\in U$,
  we have $(S+T)u=Su+Tu=0+0=0$, so $S+T\in\mathcal{E}$. Similarly,
  $(\alpha T)u=\alpha(Tu)=\alpha\cdot0=0$, so $\alpha T\in\mathcal{E}$.
  Therefore, $\mathcal{E}$ is a subspace of $\Lin(V,W)$.
\item Consider the linear map $\Phi\colon\Lin(V,W)\to\Lin(U,W)$ defined by
  $\Phi(T)=T|_U$. Let $T\in\Lin(V,W)$ such that $\Phi(T)=T|_U=0$. Then
  $T|_U=0$ implies that for every $u\in U$, we have $Tu=0$, i.e.,
  $U\subseteq\nul T$. Hence, $T\in\mathcal{E}$, showing that
  $\nul\Phi\subseteq\mathcal{E}$. Conversely, if $T\in\mathcal{E}$, then
  $U\subseteq\nul T$, which means that $T|_U=0$, i.e., $T\in\nul\Phi$.
  Therefore, $\nul\Phi=\mathcal{E}$, and $\dim\mathcal{E}=\dim\nul\Phi$.

  We now show that $\Phi$ is surjective. Let $S\in\Lin(U,W)$, and let
  $u_1,\dots,u_m$ be a basis for $U$. Extend this to a basis
  $u_1,\dots,u_m,v_1,\dots,v_n$ for $V$. Define a linear map $T\in\Lin(V,W)$ by
  setting $Tu_i=Su_i$ for $i=1,\dots,m$ and $Tv_j=0$ for $j=1,\dots,n$. Then
  $T|_U=S$, so $\Phi(T)=S$. Hence, $\Phi$ is surjective, and
  $\range\Phi=\Lin(U,W)$. Therefore,
  $\dim\range\Phi=\dim\Lin(U,W)=(\dim U)(\dim W)$. By the fundamental theorem
  of linear maps,
  \[ \dim\Lin(V,W) = \dim\nul\Phi+\dim\range\Phi
     = \dim\mathcal{E}+(\dim U)(\dim W), \]
  Observing that $\dim\Lin(V,W)=(\dim V)(\dim W)$, we get
  \[ \dim\mathcal{E} = (\dim V)(\dim W)-(\dim U)(\dim W)
     = (\dim V-\dim U)(\dim W). \qedhere \]
\end{parts}
\end{solution}

\begin{exercise}
Suppose $V$ is finite-dimensional and $S,T\in\Lin(V)$. Prove that
\[ ST \text{ is invertible} \iff S \text{ and } T \text{ are invertible}. \]
\end{exercise}
\begin{solution}
Suppose $ST$ is invertible. Then $\nul ST=\{0\}$. If $v\in\nul T$, then
$STv=S0=0$, and thus $v\in\nul ST$. Since $\nul ST=\{0\}$, we have $v=0$, and
thus $\nul T=\{0\}$. Hence, $T$ is injective, and since $V$ is
finite-dimensional, $T$ is invertible. Then $S=(ST)T^{-1}$ is a product of
invertible linear maps, and thus $S$ is invertible by Exercise~3D2.

Conversely, suppose $S$ and $T$ are invertible. Then $ST$ is a product of
invertible linear maps, and thus $ST$ is invertible by Exercise~3D2.
\end{solution}

\begin{exercise}
Suppose $V$ is finite-dimensional and $S,T,U\in\Lin(V)$ and $STU=I$. Show that
$T$ is invertible and that $T^{-1}=US$.
\end{exercise}
\begin{solution}
Since $STU=I$, then we have that $S(TU)=I$. By Exercise~3B20, then $S$ is
surjective, hence invertible with $S^{-1}=TU$. Similarly, $(ST)U=I$. By
Exercise~3B19, then $U$ is injective, hence invertible with $U^{-1}=ST$.
\[ T(US) = (TU)S = S^{-1}S = I, \quad (US)T = U(ST) = UU^{-1} = I. \]
Hence, $T$ is invertible with $T^{-1}=US$.
\end{solution}

\begin{exercise}
Show that the result in Exercise~12 can fail without the hypothesis that $V$ is
finite-dimensional.
\end{exercise}
\begin{solution}
If $V$ were not finite-dimensional, let $v_1,v_2,\dots$ be an infinite basis for
$V$. Let $S,T,U\in\Lin(V)$ be defined as follows: let $S=I$,
\[ Tv_1 = 0, \quad Tv_k = v_{k-1} \text{ for } k=2,3,\dots \]
and
\[ Uv_k = v_{k+1} \text{ for } k=1,2,\dots \]
It is clear that $TUv_k=v_k$ for all $k=1,2,\dots$, so $STU=I$. However, $T$ is
not injective since $\nul T=\spn v_1$, and thus $T$ is not invertible.
\end{solution}

\begin{exercise}
Prove or give a counterexample: If $V$ is a finite-dimensional vector space and
$R,S,T\in\Lin(V)$ are such that $RST$ is surjective, then $S$ is injective.
\end{exercise}
\begin{solution}
Since $RST$ is surjective, and $V$ is finite dimensional, it follows that $RST$
is injective, hence invertible. By Exercise~3D11, then $RST=(RS)T$ is
invertible if and only if $RS$ and $T$ are invertible. In particular, $RS$ is
invertible, and applying Exercise~3D11 again, we have that $R$ and $S$ are
invertible. Hence, $S$ is injective.
\end{solution}

\begin{exercise}
Suppose $T\in\Lin(V)$ and $v_1,\dots,v_m$ is a list in $V$ such that
$Tv_1,\dots,Tv_m$ spans $V$. Prove that $v_1,\dots,v_m$ spans $V$.
\end{exercise}
\begin{solution}
Since $Tv_1,\dots,Tv_m$ spans $V$, we have $T$ surjective. Moreover, $V$ is
finite-dimensional, so a surjective linear map on $V$ is also injective, hence
invertible. Therefore, for every $v\in V$, there exist scalars
$\alpha_1,\dots,\alpha_m$ such that
\[ v = \alpha_1Tv_1+\cdots+\alpha_mTv_m. \]
Applying $T^{-1}$ to both sides, we get
\[ T^{-1}v = \alpha_1v_1+\cdots+\alpha_mv_m. \]
Since $T^{-1}$ is surjective, every vector $w\in V$ can be written as $T^{-1}v$
for some $v\in V$, and thus $v_1,\dots,v_m$ spans $V$.
\end{solution}

\begin{exercise}
Prove that every linear map from $\F^{n,1}$ to $\F^{m,1}$ is given by a matrix
multiplication. In other words, prove that if $T\in\Lin(\F^{n,1},\F^{m,1})$,
then there exists an $m$-by-$n$ matrix $A$ such that $Tx=Ax$ for every
$x\in\F^{n,1}$.
\end{exercise}
\begin{solution}
Let $e_1,\dots,e_n$ be the standard basis of $\F^{n,1}$. Define the $k$-th
column of $A$ to be $Te_k\in\F^{m,1}$ for $k=1,\dots,n$. Then for every
$x\in\F^{n,1}$, there exist scalars $\alpha_1,\dots,\alpha_n$ such that
\[ x = \alpha_1e_1+\cdots+\alpha_ne_n. \]
Applying $T$ to both sides, we have
\[ Tx = \alpha_1Te_1+\cdots+\alpha_nTe_n = Ax, \]
where the last equality follows from the definition of matrix multiplication.
Hence, $Tx=Ax$ for every $x\in\F^{n,1}$.
\end{solution}

\begin{exercise}
Suppose $V$ is finite-dimensional and $S\in\Lin(V)$. Define
$\mathcal{A}\in\Lin\bigl(\Lin(V)\bigr)$ by
\[ \mathcal{A}(T) = ST \]
for $T\in\Lin(V)$.
\begin{parts}
\item Show that $\dim\nul\mathcal{A}=(\dim V)(\dim\nul S)$.
\item Show that $\dim\range\mathcal{A}=(\dim V)(\dim\range S)$.
\end{parts}
\end{exercise}
\begin{solution}\leavevmode
\begin{parts}
\item Consider $T\in\nul\mathcal{A}$. Then $\mathcal{A}(T)=ST=0$, so $T$ is a
  linear map from $V$ to $\nul S$. Conversely, if $T$ is a linear map from $V$
  to $\nul S$, then $ST=0$, so $T\in\nul\mathcal{A}$. Hence, $\nul\mathcal{A}$
  is the set of linear maps from $V$ to $\nul S$, and thus
  $\dim\nul\mathcal{A}=(\dim V)(\dim\nul S)$.
\item By fundamental theorem of linear maps, we have
  \[ \dim\Lin(V) = \dim\nul\mathcal{A}+\dim\range\mathcal{A}. \]
  Since $\dim\Lin(V)=(\dim V)^2$ and
  $\dim\nul\mathcal{A}=(\dim V)(\dim\nul S)$, we get
  \[ (\dim V)^2 = (\dim V)(\dim\nul S)+\dim\range\mathcal{A}, \]
  which implies
  \[ \dim\range\mathcal{A} = (\dim V)^2-(\dim V)(\dim\nul S)
     = (\dim V)(\dim\range S). \qedhere \]
\end{parts}
\end{solution}

\begin{exercise}
Show that $V$ and $\Lin(\F,V)$ are isomorphic vector spaces.
\end{exercise}
\begin{solution}
Let $T\colon\Lin(\F,V)\to V$ be defined by $T(\alpha)=\alpha(1)$ for
$\alpha\in\Lin(\F,V)$. We first show that $T$ is a linear map. Let
$\alpha,\beta\in\Lin(\F,V)$ and $\lambda\in\F$. Then
\[ T(\alpha+\beta) = (\alpha+\beta)(1) = \alpha(1)+\beta(1)
   = T\alpha+T\beta, \]
and
\[ T(\lambda\alpha) = (\lambda\alpha)(1) = \lambda\alpha(1) = \lambda T\alpha. \]
We now show that $T$ is an isomorphism. To see injectivity, let
$\alpha,\beta\in\Lin(\F,V)$ such that $T\alpha=T\beta$. Then
$\alpha(1)=\beta(1)$. For any $x\in\F$, we have
$\alpha(x)=x\alpha(1)=x\beta(1)=\beta(x)$, so $\alpha=\beta$. Hence, $T$ is
injective.

To see surjectivity, let $v\in V$. Define $\alpha\in\Lin(\F,V)$ by
$\alpha(x)=xv$ for $x\in\F$. Then $T\alpha=\alpha(1)=v$, so $T$ is surjective.
Therefore, $T$ is an isomorphism, and thus $V$ and $\Lin(\F,V)$ are isomorphic
vector spaces.
\end{solution}

\begin{exercise}
Suppose $V$ is finite-dimensional and $T\in\Lin(V)$. Prove that $T$ has the
same matrix with respect to every basis of $V$ if and only if $T$ is a scalar
multiple of the identity operator.
\end{exercise}
\begin{solution}
\solnote{Manual Remark}{See the alternate proofs at the end of this solution
for other constructions of bases to eliminate off-diagonal entries.}

Suppose $T$ has the same matrix with respect to every basis of $V$. Then for any
set of bases for $V$, its matrix must be the same. If $\dim V\le1$, then every
operator on $V$ is a scalar multiple of the identity operator. Now suppose
$\dim V\ge2$. For any bases $v_1,\dots,v_n$ and $u_1,\dots,u_n$ of $V$, we have
\[ \Mat\bigl(T,(v_1,\dots,v_n)\bigr) = \Mat\bigl(T,(u_1,\dots,u_n)\bigr). \]
Hence, we can declare that their matrix is just some $A$. We now consider the
basis $v_1+v_2,v_2,\dots,v_n$, which yields the same matrix $A$. Let
$w_1=v_1+v_2$ and $w_i=v_i$ for $i\ge2$. In the new basis, we have the
expression
\begin{align*}
Tw_1 &= A_{1,1}w_1+A_{2,1}w_2+\cdots+A_{n,1}w_n\\
     &= A_{1,1}v_1+(A_{1,1}+A_{2,1})v_2+A_{3,1}v_3+\cdots+A_{n,1}v_n.
\end{align*}
This must equal the expression for $Tw_1=Tv_1+Tv_2$ in the original basis:
\[ Tv_1+Tv_2 = (A_{1,1}+A_{1,2})v_1+(A_{2,1}+A_{2,2})v_2+\cdots
   +(A_{n,1}+A_{n,2})v_n. \]
Comparing coefficients of $v_1,\dots,v_n$ in the two expressions, we get
\begin{align*}
A_{1,1} &= A_{1,1}+A_{1,2},\\
A_{1,1}+A_{2,1} &= A_{2,1}+A_{2,2},\\
A_{j,1} &= A_{j,1}+A_{j,2} \text{ for } j\ge3.
\end{align*}
From this, we obtain that $A_{1,2}=0$, $A_{2,2}=A_{1,1}$, and $A_{j,2}=0$ for
$j\ge3$, i.e., the diagonal elements $A_{1,1}$ and $A_{2,2}$ are equal, while
off-diagonal entries in column $2$ are $0$. Now consider the basis
$v_1,v_1+v_2,\dots,v_n$. Let $u_1=v_1$, $u_2=v_1+v_2$, and $u_i=v_i$ for
$i\ge3$. In the new basis, we have
\begin{align*}
Tu_2 &= A_{1,2}u_1+\cdots+A_{n,2}u_n\\
     &= A_{1,2}v_1+A_{2,2}(v_1+v_2)+A_{3,2}v_3+\cdots+A_{n,2}v_n\\
     &= (A_{1,2}+A_{2,2})v_1+A_{2,2}v_2+A_{3,2}v_3+\cdots+A_{n,2}v_n.\\
     &= A_{1,1}v_1+A_{2,2}v_2.
\end{align*}
This must equal the expression for $Tu_2=Tv_1+Tv_2$ in the original basis,
whose expression is listed above. Comparing coefficients of $v_1$ and $v_2$, we
get
\begin{align*}
A_{1,1} &= A_{1,1}+A_{1,2},\\
A_{2,2} &= A_{2,1}+A_{2,2},\\
A_{j,2} &= A_{j,1}+A_{j,2} \quad\text{for } j\ge3.
\end{align*}
From this, we obtain similar results about diagonal entries, but the
off-diagonal entries in column $1$ must also be $0$. We may repeat the first
process for $v_1+v_k$ for $k\ge3$ to get that entries off the diagonal in column
$k$ must also be $0$. Continuing this process, we conclude all off-diagonal
entries are $0$, and all diagonal entries are equal, i.e., $A$ is a scalar
multiple of the identity matrix.

Conversely, suppose $T$ is a scalar multiple of the identity operator. Then
there exists some $\lambda\in\F$ such that $T=\lambda I$. For any basis of $V$,
the matrix representation of $T$ with respect to that basis is the diagonal
matrix with $\lambda$ on the diagonal. Hence, $T$ has the same matrix with
respect to every basis of $V$.

\emph{Alternate proof.} The manual solution provides a constructive proof
utilizing sums of the form $v_1+v_k$ to eliminate the $k$-th column, along with
$v_2+v_1$ to eliminate the 1st column. Here, we present different
constructions. Each proof assumes the declaration of the matrix to be $A$ and
that $\dim V\ge2$.

This method shows that we can swap two \emph{arbitrary} basis vectors to force
the same value of all off-diagonal entries as well as the same value of all
diagonal entries. Then, we can negate one basis vector to force all
off-diagonal entries to be $0$ via the linear independence of the basis
vectors.

\emph{Coolempire93.} Let $i,j$ be arbitrary with $1\le i,j\le n$ and $i\ne j$.
Consider the basis $w_1,\dots,w_n$ of $V$ defined by $w_i=v_j$, $w_j=v_i$ and
$w_k=v_k$ otherwise. Then
\begin{align*}
T(w_j-v_i) &= Tw_j-Tv_i\\
T(0) &= A_{1,j}w_1+A_{2,j}w_2+\cdots+A_{n,j}w_n
        -(A_{1,i}v_1+A_{2,i}v_2+\cdots+A_{n,i}v_n)\\
0 &= (A_{j,j}-A_{i,i})v_i+(A_{i,j}-A_{j,i})v_j
     +\sum_{\substack{k\ne i\\ k\ne j}}(A_{k,j}-A_{k,i})v_k.
\end{align*}
By linear independence of $v_1,\dots,v_n$, we have $A_{i,i}=A_{j,j}$,
$A_{i,j}=A_{j,i}$ and $A_{k,j}=A_{k,i}$ for all $k\ne i,j$. Because $i$ and $j$
were arbitrary, we have $A_{1,1}=A_{2,2}=\cdots=A_{n,n}=\lambda$, as well as $A$
is symmetric, and all off-diagonal entries match across rows, so we can say
$A_{i,j}=\gamma$ for all $i\ne j$. To finish, consider the basis
$u_1,\dots,u_n$ defined by $u_1=-v_1$ and $u_k=v_k$ otherwise. Then
\begin{align*}
T(v_1+u_1) &= Tv_1+Tu_1\\
T(0) &= \lambda v_1+\gamma(v_2+\cdots+v_n)
        +\bigl(\lambda u_1+\gamma(u_2+\cdots+u_n)\bigr)\\
0 &= 2\gamma(v_2+\cdots+v_n).
\end{align*}
Again by linear independence of $v_1,\dots,v_n$, we have $2\gamma=0$ and thus
$\gamma=0$. Therefore, $A=\lambda I$ and $T$ is a scalar multiple of the
identity.

\emph{Alternate proof.} This method shows that we can negate one
\emph{arbitrary} basis vector to force
all off-diagonal entries to be $0$ via the linear independence of the basis
vectors. Then, we can swap $v_1$ with any other basis vector to force all
diagonal entries to be the same.

\emph{Coolempire93.} Let $j$ be arbitrary and consider the basis
$w_1,\dots,w_n$ of $V$ defined by $w_j=-v_j$ and $w_i=v_i$ otherwise. Then
\begin{align*}
T(w_j+v_j) &= Tw_j+Tv_j\\
T(0) &= A_{1,j}w_1+A_{2,j}w_2+\cdots+A_{n,j}w_n
        +A_{1,j}v_1+A_{2,j}v_2+\cdots+A_{n,j}v_n\\
0 &= 2\sum_{k\ne j}A_{k,j}v_k.
\end{align*}
By linear independence of $v_1,\dots,v_n$, we have $A_{k,j}=0$ for all $k\ne j$.
Since $j$ was arbitrary, we have that all off-diagonal entries of $A$ are $0$.

We turn to one more basis to prove that all diagonal entries of $A$ must be the
same. Once again, let $j$ be arbitrary, and this time consider the basis
$u_1,\dots,u_n$ defined by $u_1=v_j$, $u_j=v_1$ and $u_i=v_i$ otherwise. Then
\begin{align*}
T(v_1-u_j) &= Tv_1-Tu_j\\
T(0) &= A_{1,1}v_1-A_{j,j}u_j\\
0 &= (A_{1,1}-A_{j,j})v_1.
\end{align*}
As $v_1$ belongs to a basis for $V$, it is nonzero, and thus as $j$ was
arbitrary, $A_{1,1}=A_{2,2}=\cdots=A_{n,n}=\lambda$ for some scalar $\lambda$,
and we have $A=\lambda I$, i.e., $T$ is a multiple of the identity operator.
\end{solution}

\begin{exercise}
Suppose $q\in\Poly(\R)$. Prove that there exists a polynomial $p\in\Poly(\R)$
such that
\[ q(x) = (x^2+x)p''(x)+2xp'(x)+p(3) \]
for all $x\in\R$.
\end{exercise}
\begin{solution}
Let $T\colon\Poly_n(\R)\to\Poly_n(\R)$ be defined by
\[ Tp = (x^2+x)p''+2xp'+p(3) \]
for $p\in\Poly_n(\R)$. We first show that $T$ is a linear map. Let
$p,q\in\Poly_n(\R)$ and $\alpha\in\R$. Then
\[ T(p+q) = (x^2+x)(p+q)''+2x(p+q)'+(p+q)(3) = Tp+Tq, \]
and
\[ T(\alpha p) = (x^2+x)(\alpha p)''+2x(\alpha p)'+(\alpha p)(3)
   = \alpha Tp. \]
Then $T\in\Lin\bigl(\Poly_n(\R)\bigr)$. We now show that $T$ is invertible. To
see injectivity, let $p\in\Poly_n(\R)$ such that $Tp=0$. Then
\[ (x^2+x)p''+2xp'+p(3) = 0. \]
Then
\[ (x^2+x)p''+2xp' = -p(3). \]
We analyze the degree of $p$ as follows:
\begin{itemize}
\item If $\deg p\ge2$, consider $(x^2+x)p''$. In particular,
  $(a_mx^m)''=a_mm(m-1)x^{m-2}$, hence the coefficient of $x^m$ is
  $a_mm(m-1)$. Moreover, $2xp'$ has coefficient $2ma_m$. Since $\deg p\ge2$,
  then $m\ge2$, so the total coefficient of $x^m$ is
  $a_mm(m-1)+2ma_m=a_mm(m-1+2)=a_mm(m+1)$. Since this is nonzero for
  $a_m\ne0$, we get a contradiction with the right-hand side being a constant.
\item If $\deg p=1$, then $p(x)=ax+b$ for some $a,b\in\R$. Then $p''(x)=0$ and
  $p'(x)=a$, so the left-hand side is $2ax$. Since the right-hand side is a
  constant, this is a contradiction.
\item If $\deg p=0$, then $p(x)=c$ for some $c\in\R$. Then $p''(x)=0$ and
  $p'(x)=0$, so the left-hand side is $0$, and the right-hand side is
  $p(3)=c$. Hence, $c=0$, and thus $p=0$. Therefore, $T$ is injective.
\end{itemize}
Since $T$ is injective and $\Poly_n(\R)$ is finite-dimensional, then $T$ is
surjective, hence invertible. Then there does exist some $p\in\Poly_n(\R)$ such
that $Tp=q$.
\end{solution}

\begin{exercise}
Suppose $n$ is a positive integer and $A_{j,k}\in\F$ for all
$j,k=1,\dots,n$. Prove that the following are equivalent (note that in both
parts below, the number of equations equals the number of variables).
\begin{parts}
\item The trivial solution $x_1=\cdots=x_n=0$ is the only solution to the
  homogeneous system of equations
  \[ \begin{gathered}
     \sum_{k=1}^n A_{1,k}x_k = 0\\ \vdots\\ \sum_{k=1}^n A_{n,k}x_k = 0.
     \end{gathered} \]
\item For every $c_1,\dots,c_n\in\F$, there exists a solution to the system of
  equations
  \[ \begin{gathered}
     \sum_{k=1}^n A_{1,k}x_k = c_1\\ \vdots\\ \sum_{k=1}^n A_{n,k}x_k = c_n.
     \end{gathered} \]
\end{parts}
\end{exercise}
\begin{solution}\leavevmode
\begin{parts}
\item Suppose the trivial solution $0\in\F^{n,1}$ is the only solution to the
  homogeneous system of equations. Then $\nul T=\{0\}$, where
  $T\colon\F^{n,1}\to\F^{n,1}$ is defined by $Tx=Ax$ for $x\in\F^{n,1}$. Since
  $\nul T=\{0\}$, then $T$ is injective. Since $\F^{n,1}$ is finite-dimensional,
  then $T$ is surjective, hence invertible. Then for every
  $c_1,\dots,c_n\in\F$, there exists some $x\in\F^{n,1}$ such that $Tx=c$,
  where $c$ is the column vector with entries $c_1,\dots,c_n$. Hence, there
  exists a solution to the system of equations.
\item Suppose for every $c_1,\dots,c_n\in\F$, there exists a solution to the
  system of equations. Then $T\colon\F^{n,1}\to\F^{n,1}$ defined by $Tx=Ax$ for
  $x\in\F^{n,1}$ is surjective. Since $\F^{n,1}$ is finite-dimensional, then
  $T$ is injective. Hence, the trivial solution is the only solution to the
  homogeneous system of equations.
\end{parts}
\end{solution}

\begin{exercise}
Suppose $T\in\Lin(V)$ and $v_1,\dots,v_n$ is a basis of $V$. Prove that
\[ \Mat\bigl(T,(v_1,\dots,v_n)\bigr) \text{ is invertible} \iff
   T \text{ is invertible}. \]
\end{exercise}
\begin{solution}
Suppose $\Mat\bigl(T,(v_1,\dots,v_n)\bigr)$ is invertible. Then there exists
some $A\in\F^{n,n}$ such that $A\Mat\bigl(T,(v_1,\dots,v_n)\bigr)=I$. Let $U$
be the linear map from $V$ to $V$ whose matrix with respect to the basis
$v_1,\dots,v_n$ is $A$. Then
\begin{align*}
\Mat\bigl(UT,(v_1,\dots,v_n)\bigr)
  &= \Mat\bigl(U,(v_1,\dots,v_n)\bigr)\Mat\bigl(T,(v_1,\dots,v_n)\bigr)\\
  &= A\Mat\bigl(T,(v_1,\dots,v_n)\bigr) = I.
\end{align*}
Then $UT$ is the identity operator, and thus $T$ is injective. Since $V$ is
finite-dimensional, then $T$ is surjective, hence invertible.

Suppose $T$ is invertible. Then $T^{-1}$ exists, and we have
\[ \Mat\bigl(T^{-1}T,(v_1,\dots,v_n)\bigr)
   = \Mat\bigl(T^{-1},(v_1,\dots,v_n)\bigr)\Mat\bigl(T,(v_1,\dots,v_n)\bigr)
   = I. \]
Hence, $\Mat\bigl(T,(v_1,\dots,v_n)\bigr)$ is invertible.
\end{solution}

\begin{exercise}
Suppose that $u_1,\dots,u_n$ and $v_1,\dots,v_n$ are bases of $V$. Let
$T\in\Lin(V)$ be such that $Tv_k=u_k$ for each $k=1,\dots,n$. Prove that
\[ \Mat\bigl(T,(v_1,\dots,v_n)\bigr)
   = \Mat\bigl(I,(u_1,\dots,u_n),(v_1,\dots,v_n)\bigr). \]
\end{exercise}
\begin{solution}
By definition, column $k$ of $\Mat\bigl(T,(v_1,\dots,v_n)\bigr)$ is $Tv_k=u_k$
with respect to the basis $v_1,\dots,v_n$. Similarly, column $k$ of
$\Mat\bigl(I,(u_1,\dots,u_n),(v_1,\dots,v_n)\bigr)$ is also $u_k$ with respect
to the basis $v_1,\dots,v_n$.
\end{solution}

\begin{exercise}
Suppose $A$ and $B$ are square matrices of the same size and $AB=I$. Prove that
$BA=I$.
\end{exercise}
\begin{solution}
Let $T\colon\F^{n,1}\to\F^{n,1}$ be defined by $Tx=Ax$ for $x\in\F^{n,1}$.
Since $AB=I$, then $T$ is surjective. Since $\F^{n,1}$ is finite-dimensional,
surjectivity implies injectivity. Hence, $T$ is invertible, and
$T(Bx)=x=T(T^{-1}x)$. Then $Bx=T^{-1}x$ for every $x\in\F^{n,1}$. Therefore,
\[ BAx = B(Tx) = T^{-1}(Tx) = x \]
for every $x\in\F^{n,1}$. Hence, $BA=I$.
\end{solution}

\vspace{2em}
\end{exc}
